\chapter{Projekt bazy danych}
\label{ch:baza}

Baza danych systemu sk\l{}ada si\k{e} z~pi\k{e}ciu tabel: \texttt{Kursanci},
\texttt{Instruktorzy}, \texttt{Pojazdy}, \texttt{RezerwacjeJazd}
i~\texttt{Administratorzy}. Tabela \texttt{RezerwacjeJazd} jest tabel\k{a}
centraln\k{a} -- {\l}\k{a}czy ze sob\k{a} pozosta\l{}e trzy przez powi\k{a}zania.

\section{Diagram ERD}
\label{sec:erd}

Poni\.{z}szy diagram przedstawia struktur\k{e} bazy danych oraz powi\k{a}zania
mi\k{e}dzy tabelami.

\begin{figure}[H]
  \centering
  \includegraphics[width=\textwidth]{img/erd}
  \caption{Diagram ERD bazy danych systemu OSK}
  \label{fig:erd}
\end{figure}

\section{Normalizacja}

Schemat bazy danych spe\l{}nia trzeci\k{a} posta\'{c} normaln\k{a} (3NF).
Ka\.{z}da tabela przechowuje dane dotycz\k{a}ce tylko jednego obiektu:
osobno kursanci, osobno instruktorzy, osobno pojazdy. Powi\k{a}zania
mi\k{e}dzy nimi realizowane s\k{a} przez klucze obce w~tabeli rezerwacji,
bez powt\'{o}rzenia danych.

\section{Relacje i integralno\'{s}\'{c} danych}

Ka\.{z}da rezerwacja musi by\'{c} powi\k{a}zana z~istniej\k{a}cym kursantem,
instruktorem i~pojazdem. System blokuje usuni\k{e}cie kursanta, instruktora
lub pojazdu, je\.{z}eli ma on przypisane jazdy w~systemie -- zapobiega to
przypadkowej utracie danych.

Loginy u\.{z}ytkownik\'{o}w oraz numer PESEL kursanta musz\k{a} by\'{c}
unikalne. Wymuszane jest to na poziomie bazy danych, co uniemo\.{z}liwia
duplikaty nawet przy jednoczesnym dost\k{e}pie.

Struktura bazy zosta\l{}a utworzona przy u\.{z}yciu mechanizmu migracji --
zmiany schematu s\k{a} wersjonowane i~mo\.{z}na je w~ka\.{z}dej chwili
odtworzy\'{c}.
